% TOP_PROBLEM: {"id": 2306813, "title": "Research Problems in Function Theory — Problem 6.13′", "category_id": 8, "category": {"id": 8, "name": "analysis", "display_name": "Analysis", "description": "Limits, continuity, calculus, and function theory.", "slug": "analysis", "order_index": 8, "created_at": "2026-07-31T15:26:25.671Z"}, "status": "open"}
% FIRST_POSTED: null
% ENTRY_KIND: statement_only
% Imported from: batch04.tex; UnsolvedMath ID 2306813
\documentclass[11pt]{article}
\usepackage[margin=25mm]{geometry}
\usepackage{fontspec}
\setmainfont{Latin Modern Roman}
\usepackage{amsmath,amssymb,mathtools}
\usepackage{unicode-math}
\setmathfont{Latin Modern Math}
\usepackage{xurl,hyperref}
\hypersetup{colorlinks=true,urlcolor=blue}
\setcounter{secnumdepth}{0}
\setlength{\parindent}{0pt}
\setlength{\parskip}{5pt}
\setlength{\emergencystretch}{3em}
\newcommand{\sref}[2]{\hyperlink{source:#1:#2}{[S#2]}}
\newcommand{\cataloguescope}[1]{\paragraph{Recorded target.}#1}
\begin{document}
% UnsolvedMath ID 2306813; source catalogue code AMR-022-6813
\setcounter{section}{2306812}
\section{Research Problems in Function Theory — Problem 6.13$\prime$}\label{2306813}
{\small\sffamily UnsolvedMath ID 2306813 · Analysis}
\subsection{Definitions and mathematical statement}

\paragraph{Shared notation.}$\mathbb N=\{1,2,\ldots\}$, and $\mathbb Z,\mathbb Q,\mathbb R,\mathbb C$ denote the integers, rationals, reals and complex numbers. Any other conventions are specified locally.


Write $\mathbb D=\{z\in\mathbb C:|z|<1\}$ and $\mathbb T=\partial\mathbb D$. For $p>0$ and a holomorphic $f$ on $\mathbb D$, let $n_f(w)$ count the preimages of $w$ with multiplicity. The function is areally mean $p$-valent when
\[\int_0^{2\pi}\int_0^R n_f(\rho e^{i\theta})\rho\,d\rho\,d\theta\leq\pi pR^2\qquad(R>0).\]
Writing $f(z)=\sum_{j=0}^{\infty}a_jz^j$, put $\mu_p(f)=\max\{|a_j|:j\in\mathbb Z,\ 0\leq j\leq p\}$, and let $S(p)$ be the mean $p$-valent functions with $\mu_p(f)=1$.
Fix $p>1/4$ and an integer $m\geq1$. Call an exponent $a\in\mathbb R$ admissible if for every mean $p$-valent $f(z)=\sum_{j\geq0}a_jz^j$ and every integer $k\geq1$ there is $C=C(f,k,m,p)<\infty$ such that, for each sufficiently large integer $n$, one can choose complex numbers $c_0,\ldots,c_m$ with $c_0=1$, $|c_m|\geq4^{-m}$, and
\[\left|\sum_{\nu=0}^{m}c_\nu a_{n+j+\nu}\right|\leq Cn^a\qquad(0\leq j\leq k).\]
The coefficients $c_\nu$ may depend on $f,k,m,n$. Let $\alpha_m(p)$ be the infimum of admissible exponents. This definition applies to the full mean-valent class, without the normalization $\mu_p=1$; the bound may depend on $f$.
\paragraph{Statement recorded in the supplied source.}
Determine the sharp exponents $\alpha_m(p)$ and whether each boundary exponent is admissible. In particular, are
\[\alpha_1(p)=p-1\quad(1/2<p<1),\qquad
\alpha_1(p)=-1/2\quad(1/4<p<1/2),\]
and $\alpha_m(p)=-1/2$ whenever $m+1>4p$? The source reports the available admissible exponent $-1/2+8p^{3/2}/\sqrt m$ and the obstruction $\alpha_m(p)\geq2p/(m+1)-1$. \sref{2306813}{2}
\subsection{Short English statement}
Determine the optimal powers controlling simultaneous finite linear combinations of consecutive coefficients, including the stated mean-valence conjectures.
\subsection{Sources}
\begin{description}
\item[\hypertarget{source:2306813:1}{S1}] ULAM.AI and the UnsolvedMath Contributors, UnsolvedMath: A Curated Collection of Open Mathematics Problems, record 2306813 (AMR-022-6813). CC BY 4.0. \url{https://huggingface.co/datasets/ulamai/UnsolvedMath}.
\item[\hypertarget{source:2306813:2}{S2}] W. K. Hayman and E. F. Lingham, Research Problems in Function Theory (2018), Problem 6.13 prime, complete problem and update, Problem6.13, equation(6.5), complete mean-p notation, examples and update. \url{https://arxiv.org/abs/1809.07200}.
\end{description}
\label{end:2306813}
\end{document}
